\documentclass[letterpaper, 12pt, titlepage]{article} % Tipo de documento

% ================================================================
% PAQUETES BÁSICOS Y CONFIGURACIÓN ESENCIAL
% ================================================================
\usepackage[utf8]{inputenc}         % Codificación de caracteres UTF-8
\usepackage[spanish]{babel}         % Localización en español (guiones, títulos)
\usepackage[margin=1in]{geometry}   % Configuración de márgenes de página

% ================================================================
% TIPOGRAFÍA Y FORMATO DE TEXTO
% ================================================================
\usepackage{helvet}                         % Fuente Sans-Serif (Helvetica-like)
\renewcommand{\familydefault}{\sfdefault}   % Establece Helvet como fuente principal
\usepackage{setspace}                       % Control de interlineado (1.5, doble espacio)
\usepackage{parskip}                        % Espaciado entre párrafos (en lugar de sangría)
\usepackage{ragged2e}                       % Mejor manejo de justificación de texto

% ================================================================
% ELEMENTOS GRÁFICOS Y DISEÑO
% ================================================================
\usepackage{graphicx}               % Manejo de imágenes (inserción y escalado)
\usepackage{xcolor}                 % Sistema de colores personalizados
\usepackage{pdflscape}              % Páginas en orientación horizontal

% ================================================================
% TABLAS PROFESIONALES
% ================================================================
\usepackage{booktabs}              % Reglas tipográficas para tablas (formato APA)
\usepackage{xltabular}             % Tablas largas que ocupan todo el ancho del texto
% ================================================================
% LISTAS Y ESTRUCTURAS
% ================================================================
\usepackage{enumitem}                   % Personalización avanzada de listas
\setlist{nosep, leftmargin=*, wide=0pt} % Listas ultra compactas

% ================================================================
% HIPERVÍNCULOS Y METADATOS PDF
% ================================================================
\usepackage{hyperref}              % Enlaces interactivos y metadatos PDF
\hypersetup{
    colorlinks = true,             % Colores en lugar de cajas
    linkcolor  = blue!70,          % Color para enlaces internos
    urlcolor   = cyan!80,          % Color para URLs
    citecolor  = magenta!70,       % Color para citas
}

% ================================================================
% CONFIGURACIONES FINALES
% ================================================================
\renewcommand{\thesection}{\Roman{section}} % Numeración romana para secciones
\newcommand{\universidad}{Universidad de las Regiones Autónomas de la Costa Caribe Nicaragüense} % Nombre de la universidad
\newcommand{\carrera}{Nombre de la Carrera} % Nombre de la carrera
\newcommand{\asignatura}{Nombre de la Asignatura} % Nombre de la asignatura
\newcommand*{\thead}[1]{\multicolumn{1}{>{\centering\arraybackslash}X}{\bfseries #1}} % Encabezado centrado y en negritas en tablas

\begin{document}

% Portada del documento
% titlepage incluye un \pagebreak (salto de página) automático al final

\begin{titlepage}
    \centering % centrar el contenido de la portada
    % \minipage permite dividir la portada en dos columnas
    \begin{minipage}{0.25\textwidth} % logo en la columna izquierda
        \centering % centrar el logo dentro de la \minipage
        \includegraphics[width=0.6\textwidth]{images/uraccan_logo.png} % inclusión del logo de la universidad
    \end{minipage}%
    \begin{minipage}{0.75\textwidth} % texto en la columna derecha
        \centering % centrar el texto dentro de la \minipage
        \Large{\textbf{\universidad}}\\ % Nombre de la universidad, configurado en PLANTILLA-SILABO-URACCAN.tex
        \Large{\textbf{URACCAN}}\\ % Siglas de la universidad
    \end{minipage}

    \vfill % espacio flexible para centrar el contenido verticalmente
    \large{\textbf{Syllabus}}\\
    \Huge{\textbf{\asignatura}}\\ % Nombre de la asignatura, configurado en PLANTILLA-SILABO-URACCAN.tex

    \vfill % espacio flexible para centrar el contenido verticalmente
    \small{\textbf{\today{}}}\\ % Fecha actual, configurada con \today, que calcula la fecha automáticamente
\end{titlepage}
\section{Datos Generales}
% Para poder usar \footnotes en el título de la sección:
% \SECTION[Título interno de la sección]{Título visible de la sección}
\section[Descriptor de la asignatura]{Descriptor de la asignatura\footnote{Texto íntegro del descriptor, tomado del programa de la asignatura.}}

% Aquí debe poner el texto del descriptor de la asignatura, se encuentra en el programa de la asignatura.

% Texto Lorem Ipsum para el descriptor de la asignatura. Eliminar en producción.
\lipsum[1-4] % Genera cuatro párrafos de texto de ejemplo
\pagebreak
% Para poder usar footnotes en el título de la sección:
% \SECTION[Título interno de la sección]{Título visible de la sección}
\section[Objetivos Generales o Resultados de Aprendizaje]{Objetivos Generales o Resultados de Aprendizaje\footnote{Texto íntegro de los objetivos o resultados de aprendizaje, tomado del programa de la asignatura.}}
%Eliminar \PENDIENTE en producción.
\pendiente{Según el enfoque de la asignatura, se pueden definir objetivos generales o resultados de aprendizaje. Pero no los dos a la vez. Revise el programa de la asignatura.}

\small
\begin{xltabular}{\linewidth}{@{}X X X@{}}
    \toprule
    \encabezadodetabla{\textbf{Conceptuales}} & \encabezadodetabla{\textbf{Procedimentales}} & \encabezadodetabla{\textbf{Actitudinales}} \\
    \midrule
    \begin{itemize}
        \item \lipsum[1][1-3] % Texto de ejemplo, eliminar en producción
        \item \lipsum[2][1-3] % Texto de ejemplo, eliminar en producción
        \item \lipsum[3][1-3] % Texto de ejemplo, eliminar en producción
        \item \lipsum[4][1-3] % Texto de ejemplo, eliminar en producción
    \end{itemize} &
    \begin{itemize}
        \item \lipsum[1][1-3] % Texto de ejemplo, eliminar en producción
        \item \lipsum[2][1-3] % Texto de ejemplo, eliminar en producción
        \item \lipsum[3][1-3] % Texto de ejemplo, eliminar en producción
        \item \lipsum[4][1-3] % Texto de ejemplo, eliminar en producción
    \end{itemize} &
    \begin{itemize}
        \item \lipsum[1][1-3] % Texto de ejemplo, eliminar en producción
        \item \lipsum[2][1-3] % Texto de ejemplo, eliminar en producción
        \item \lipsum[3][1-3] % Texto de ejemplo, eliminar en producción
        \item \lipsum[4][1-3] % Texto de ejemplo, eliminar en producción
    \end{itemize} \\
    \bottomrule
\end{xltabular}
\pagebreak
\begin{landscape}
    \section{Contenidos}

    \doublespacing
    \begin{xltabular}{\linewidth}{@{}c l c c c c@{}}
        \toprule
        & & \multicolumn{4}{c}{\textbf{Horas}}\\
        \cmidrule(lr){3-6}
        \encabezadodetabla{No.} & \multicolumn{1}{c}{\textbf{Unidad}} & \encabezadodetabla{Teóricas} & \encabezadodetabla{Prácticas} & \encabezadodetabla{Trabajo Independiente} & \encabezadodetabla{Total} \\
        \midrule
        I & \lipsum[1][1] & 8 & 8 & 16 & 32 \\
        II & \lipsum[2][1] & 8 & 8 & 16 & 32 \\
        III & \lipsum[3][1] & 8 & 8 & 16 & 32 \\
        IV & \lipsum[4][1] & 8 & 8 & 16 & 32 \\
        V & \lipsum[5][1] & 8 & 8 & 16 & 32 \\
        VI & \lipsum[6][1] & 8 & 8 & 16 & 32 \\
        \midrule
        \multicolumn{2}{c}{\textbf{Evaluación}} & 0 & 4 & 8 & 12 \\
        \midrule
        \multicolumn{2}{c}{\textbf{TOTAL}} & 30 & 34 & 128 & 192 \\
        \bottomrule
    \end{xltabular}
\end{landscape}
\section{Estrategias Metodológicas}

Las actividades, métodos y técnicas para lograr un resultado de aprendizaje efectivo, se deben realizar usando las diferentes estrategias que se enumeran a continuación:

\textbf{Interculturalidad y género}, este curso abordará la interculturalidad tomando las cualidades siguientes:\

\begin{itemize}
    \item Promover actitudes y capacidades de aceptación, respeto e interacción en la diversidad cultural y género.
    \item Mediar conflictos a través de la escucha activa, empatía y ética para propiciar clima de convivencia y tranquilidad.
    \item Comprender al estudiantado brindándole confianza en el desarrollo de las actividades académicas que potencien su formación técnica, profesional y humanística.
    \item Garantizar las oportunidades con igualdad de condiciones al estudiantado para que demuestren sus capacidades, habilidades y destrezas acompañándolos en las experiencias didácticas.
    \item Promover estrategias de trabajo colaborativos y cooperativos que favorezcan las relaciones intrapersonales e interpersonales que contribuyan a eliminar posibles estereotipos.
    \item Promover el empoderamiento de hombres y mujeres vulnerables.
\end{itemize}

\textbf{Conferencias}, donde se pueden abordar todos los aspectos teóricos y se pueden aclarar al mismo tiempo aspectos iniciales y metodología a cumplir durante cada sesión de clases.

\textbf{Talleres}, donde se le facilite al estudiantado el camino a seguir para aplicar los conocimientos técnicos adquiridos durante la sesión teórica, donde se puede verificar si el estudiante ha asimilado cada contenido, donde el método de presentación es libre, dejando como opciones los reportes, videos y presentación en vivo.

\textbf{Tecnologías de la información y comunicaciones (TIC)}, basado en dispositivos tecnológicos e inteligentes con acceso a una amplia variedad de fuentes de información, a la vez del uso de software educativo para fortalecer la enseñanza aprendizaje con herramientas innovadoras, tales como la plataforma educativa MOODLE, Telegram como método de comunicación directa, asimismo, la plataforma de alojamiento de código fuente GitHub.

\textbf{Exposiciones}, que permitan al estudiante indagar sobre aspectos teóricos y técnicos referentes al proceso de diseño y escritura de páginas web, así como la utilización adecuada y buen manejo de las tecnologías web.

\textbf{Resolución de retos prácticos} apoyándose en la guía metodológica otorgada por el docente, donde tendrán que resolverla en un plazo de tiempo acorde a la dificultad del reto, para entregar resultado al docente, donde cada estudiante, durante la sesión de clases, analizará la guía y compartirá sus dudas.

\textbf{Retroalimentación} de cada clase anterior con el objetivo de establecer vínculos en serie entre cada una, donde este espacio se usará para enfocarse en elementos que dificultaron en la práctica.

\pagebreak
\section{Materiales y Recursos Didácticos}

\subsection*{Materiales físicos y espacios}

\begin{xltabular}{\linewidth}{@{}>{\bfseries}X X@{}}
    \toprule
    Pizarra acrílica & Superficie para escritura y explicaciones en tiempo real. \\
    \midrule
    Pantalla o proyector & Dispositivo de visualización de contenido digital/audiovisual. \\
    \midrule
    Laboratorio de computación & Espacio equipado con equipos tecnológicos y acceso a internet para ejercicios prácticos. \\
    \midrule
    Computadoras & Dispositivos para ejecutar software, acceder a recursos digitales, realización de prácticas y elaboración de informes. \\
    \midrule
    Acceso a internet & Conexión necesaria para la búsqueda de información, uso de plataformas educativas y recursos en línea. \\
    \bottomrule
\end{xltabular}

\subsection*{Documentación}

\begin{xltabular}{\linewidth}{@{}>{\bfseries}X X@{}}
    \toprule
    Sílabo de la asignatura & Documento que detalla los objetivos, contenidos y metodología del curso. \\
    \midrule
    Planes de clase & Guías que describen las actividades y recursos para cada sesión. \\
    \midrule
    Bibliografía básica y complementaria & Listado de textos y recursos recomendados para profundizar en los temas tratados. \\
    \midrule
    Materiales de lectura y recursos digitales & Artículos, videos y otros recursos en línea que apoyan el aprendizaje. \\
    \midrule
    Guías de lectoescritura & Documentos que orientan sobre la correcta interpretación y redacción de textos académicos y presentaciones. \\
    \midrule
    Guías de prácticas de laboratorio & Instrucciones detalladas para la realización de experimentos, simulaciones y prácticas en el laboratorio. \\
    \bottomrule
\end{xltabular}

\subsection*{Herramientas digitales y software}

\begin{xltabular}{\linewidth}{@{}>{\bfseries}X X@{}}
    \toprule
    Sistema operativo & Software base que gestiona los recursos del hardware y permite la ejecución de aplicaciones. \\
    \midrule
    Plataforma virtual MOODLE & Entorno de aprendizaje en línea para la gestión de cursos, recursos y actividades. \\
    \midrule
    Microsoft Word o alternativas & Procesador de texto para la elaboración de documentos, informes y trabajos escritos. \\
    \midrule
    Microsoft PowerPoint o alternativas & Herramienta para la creación de presentaciones visuales que apoyan las exposiciones orales. \\
    \midrule
    Telegram o WhatsApp & Mensajería instantánea para la comunicación entre estudiantes y docentes. \\
    \midrule
    Microsoft Teams & Plataforma de colaboración y comunicación para clases virtuales, reuniones y trabajo en grupo. \\
    \bottomrule
\end{xltabular}
\pagebreak
\section{Metodología de Evaluación}

La aprobación satisfactoria del curso depende de una nota mayor o igual a 60 puntos que resulta del promedio de dos evaluaciones parciales a lo largo del curso.

La primera prueba parcial está estructurada en 60\% de evaluación sumativa y 40\% de evaluación formativa. La evaluación formativa estará dividida en dos partes: un examen teórico a realizarse en el aula virtual; y la defensa de una práctica de laboratorio basada en el contenido práctico visto en las prácticas de laboratorio del parcial.

La segunda prueba parcial también se estructura del 60\% de evaluación sumativa y 40\% de evaluación formativa enfocado en un proyecto final donde aplique todo lo aprendido.

Las evaluaciones sumativas se componen de prácticas de laboratorio y actividades de lectoescritura que resultan en exposiciones o informes de investigación de cada tema propuesto.

Las evaluaciones formativas se enfocarán en guiar al estudiante a trabajar en un proyecto de magnitud o dificultad considerable, aplicando de forma integral los conocimientos adquiridos en el transcurso de la asignatura a un caso simulado, dando una solución satisfactoria al reto presentado.

Es recomendable que las evaluaciones prácticas se valore el trabajo individual, para garantizar el aprendizaje total en cada uno.
\pagebreak
\section{Bibliografía}
%* Aquí debe poner la bibliografía de la asignatura
%! Use el formato APA para las referencias

%! Eliminar este \pendiente en producción
\pendiente{Agregar bibliografía de la asignatura.}

\subsection*{Referencias Básicas}

%! Ejemplos de libros impresos
\hangindent=0.5in
Merino, C., \& Cañizares, R. (2011). \textit{Implementación de un Sistema de Gestión de Seguridad de la Información según ISO 27001}. Fundación Confemetal.

\hangindent=0.5in
Peltier, J. (2020). \textit{Information security policies, procedures, and standards: Guidelines for effective information security management} (3ª ed.). Auerbach Publications.

%! Ejemplos de artículos de revistas
\hangindent=0.5in
Arévalo, F. M., \& Cedillo, L. P. (2017). \textit{Metodología Ágil para la gestión de Riesgos Informáticos}. Killkana Técnica, 1(2), 45-67.

\hangindent=0.5in
ISO/IEC. (2013). \textit{ISO/IEC 27005:2013 Information technology - Security techniques - Information security risk management}. ISO Bulletin, 4(2), 12-18. \href{https://doi.org/10.1016/j.isac.2013.05.002}{https://doi.org/10.1016/j.isac.2 013.05.002}

%! Ejemplos de capítulos de libros
\hangindent=0.5in
Bautista Torres, L. A. (2018). Seguridad de la información en entornos empresariales. En J. González \& M. López (Eds.), \textit{Gestión de riesgo en la seguridad informática} (pp. 78-105). Editorial Tecnológica.

\hangindent=0.5in
Imbaquinto, Esparza, D. E., Pusda, Chulde, M. R., \& Jácome León, J. G. (2016). Fundamento de Auditoría Basados en Riesgos. En \textit{Auditoría informática en entornos digitales} (pp. 112-134). UTN.

\subsection*{Referencias Complementarias}
\pagebreak
\begin{landscape}
    \section{Plan Calendario}

    \begin{adjustbox}{width=1.3847\textwidth}
        \begin{tabular}{ | p{3cm} | p{5cm} | p{5cm} | p{3cm} | p{3cm} | p{3cm} | p{5cm} | }
            \hline
            \textbf{Fecha}            & \textbf{Objetivos o Resultados de Aprendizaje} & \textbf{Unidad y Contenidos} & \textbf{Estrategias Metodológicas} & \textbf{Trabajo Independiente} & \textbf{Mecanismos de Evaluación} & \textbf{Bibliografía} \\
            \hline
            % Inicio Encuentro 1
            % Fecha
            03 de marzo 2025          &
            % Objetivos o Resultados de Aprendizaje
            \begin{minipage}
                [t]{5cm}
                \begin{itemize}[leftmargin=10pt]
                    \item Reconoce la importancia de la ética en la vida profesional.
                    \item Identifica los principios éticos que rigen la profesión.
                \end{itemize}
                \vspace{0.2cm}
            \end{minipage} &
            \begin{minipage}
                [t]{5cm}
                \begin{itemize}[leftmargin=10pt]
                    \item Unidad 1: Ética y Profesión.
                    \item Contenidos: Concepto de ética, ética profesional, principios éticos.
                \end{itemize}
                \vspace{0.2cm}
            \end{minipage} &
            \begin{minipage}
                [t]{3cm}
                \begin{itemize}[leftmargin=10pt]
                    \item Exposición.
                    \item Debate.
                    \item Análisis de casos.
                \end{itemize}
            \end{minipage} &
            \begin{minipage}
                [t]{3cm}
                \begin{itemize}[leftmargin=10pt]
                    \item Lectura de textos.
                    \item Investigación.
                \end{itemize}
            \end{minipage} &
            \begin{minipage}
                [t]{3cm}
                \begin{itemize}[leftmargin=10pt]
                    \item Participación en clase.
                    \item Exposición.
                    \item Análisis de casos.
                \end{itemize}
            \end{minipage} &
            \begin{minipage}
                [t]{5cm}
                \begin{itemize}[leftmargin=10pt]
                    \item \textit{Ética para Amador} - Fernando Savater.
                    \item \textit{Ética Profesional} - José A. Silié Ruiz.
                \end{itemize}
            \end{minipage}                                                                                                                                                                                                        \\
            \hline
            % Fin Encuentro 1
        \end{tabular}
    \end{adjustbox}
\end{landscape}

\end{document}